\documentclass[dvisvgm,tikz,border=3pt]{standalone}
\usepackage{amsmath,amssymb}
\usepackage{pgfplots}
\usepackage{CJKutf8}
\pgfplotsset{compat=1.18}
\usetikzlibrary{arrows.meta,calc}

\definecolor{themebg}{HTML}{2e362c}
\definecolor{themefg}{HTML}{edf4e8}
\definecolor{themegray}{HTML}{9ea897}
\definecolor{themecurve}{HTML}{7eb6ff}
\definecolor{themealert}{HTML}{ff7b7b}
\definecolor{themeamber}{HTML}{f5b942}

\begin{document}
\begin{CJK*}{UTF8}{gbsn}
\pagecolor{themebg}
\begin{tikzpicture}[color=themefg, text=themefg]
\begin{axis}[
    name=mainax,
    width=12.5cm,
    height=8.0cm,
    axis lines=middle,
    axis line style={color=themefg, -{Stealth}},
    tick style={color=themefg},
    tick label style={color=themefg, font=\small},
    every axis label/.style={color=themefg},
    xmin=-4.6, xmax=4.9,
    ymin=-2.4, ymax=3.2,
    xtick={-4, -3, -2, -1, 0, 1, 2, 3, 4},
    ytick={0},
    clip=false,
    xlabel={$x$},
    ylabel={$y$},
    title style={at={(0.5,1.08)},anchor=south,color=themefg,font=\bfseries\normalsize},
    title={有界性几何模型：输出全域被困在有限水平带 $[m, M]$ 内},
]

  % 有界性母图只表达“输出被统一夹住”，不能偷带周期、奇偶或单调等额外结构。
  % 先画水平带背景，再画边界；背景层不遮挡坐标轴、刻度与函数曲线。
  \fill[blue!18!themebg, opacity=0.28] (axis cs:-4.4,-1.5) rectangle (axis cs:4.4,2.0);
  \draw[dashed, themegray, line width=1.1pt] (axis cs:-4.4, 2.0) -- (axis cs:4.4, 2.0);
  \draw[dashed, themegray, line width=1.1pt] (axis cs:-4.4, -1.5) -- (axis cs:4.4, -1.5);

  % 边界说明统一放在最右侧空白端，避免压住刻度与数据。
  \node[text=themegray, fill=themebg, inner sep=0.8pt, anchor=west] at (axis cs:4.45, 2.0) {\small 上界 $y=M$};
  \node[text=themegray, fill=themebg, inner sep=0.8pt, anchor=west] at (axis cs:4.45, -1.5) {\small 下界 $y=m$};

  % 一般化非周期、非奇偶、非单调有界曲线。
  % 有理衰减项保证全局有界；不使用正弦波，避免把“有界”误认成“周期振荡”。
  \addplot[themecurve, line width=2.2pt, domain=-4.2:4.2, samples=240]
    {0.2 + 1.8*(0.9*x/(1+x^2) + 0.45*(x+1.2)/(1+(x+1.2)^2) - 0.3*(x-2)/(1+(x-2)^2))};
  \node[text=themecurve, fill=themebg, inner sep=0.8pt, anchor=south east] at (axis cs:3.95, 0.72) {\small $y=f(x)$};

  % 结论标签放在数据稀疏区，只表达像集被包含，不冒充最值已经取到。
  \node[text=themealert, fill=themebg, inner sep=0.8pt, anchor=west] at (axis cs:2.25, -1.05) {\small $R_f\subseteq[m,M]$};

\end{axis}

% Bottom summary caption placed completely OUTSIDE the axis coordinate box
\node[font=\small\bfseries, color=themefg, anchor=north] at ($(mainax.south)-(0, 0.55cm)$) {充要条件：$\exists M_0 > 0$，使得 $\forall x \in D$，均有 $\lvert f(x) \rvert \le M_0$};

\end{tikzpicture}
\end{CJK*}
\end{document}
